% TOP_PROBLEM: {"id": 10000044, "title": "From a large percolation component to a giant component on expanders", "category_id": 19, "category": {"id": 19, "name": "probability", "display_name": "Probability", "description": "Problems involving probability theory, stochastic processes, and random structures.", "slug": "probability", "order_index": 19, "created_at": "2026-07-31T15:26:25.671Z"}, "status": "open", "problem_number": "AMR-099-0044", "set_id": 13, "statusQualification": "Makes the asymptotic family and intrinsic graph-diameter convention explicit in the source’s finite-graph formulation."}
% FIRST_POSTED: 2026-10-01
% ENTRY_KIND: statement_only
% UnsolvedMath ID 10000044; source catalogue code AMR-099-0044
% !TeX program = lualatex
% Definition and sources only; this file is not an LLM solution attempt.
\documentclass[11pt,a4paper]{article}
\usepackage[margin=25mm]{geometry}
\usepackage{fontspec}
\setmainfont{lmroman10-regular.otf}[BoldFont=lmroman10-bold.otf,
 ItalicFont=lmroman10-italic.otf,BoldItalicFont=lmroman10-bolditalic.otf]
\usepackage{amsmath,amssymb,mathtools}
\usepackage{unicode-math}
\setmathfont{latinmodern-math.otf}
\AtBeginDocument{\let\setminus\smallsetminus}
\usepackage{xurl,hyperref}
\hypersetup{colorlinks=true,urlcolor=blue}
\setcounter{secnumdepth}{0}
\setlength{\parindent}{0pt}
\setlength{\parskip}{5pt}
\setlength{\emergencystretch}{3em}
\begin{document}
\section*{From a large percolation component to a giant component on expanders}\label{10000044}
{\small UnsolvedMath ID 10000044 · AMR-099-0044 · Probability · AMR Open Problem Lists}
\subsection{Definitions and mathematical statement}
Let $(G_n)$ be finite connected graphs with $|V(G_n)|\to\infty$, maximum degrees bounded by one constant $D$, and edge expansion uniformly bounded below: some $h>0$ satisfies
\[
|\{xy\in E(G_n):x\in A,y\notin A\}|\ge h|A|
\quad(0<|A|\le |V(G_n)|/2).
\]
Perform independent bond percolation with parameter $1/2$ on each $G_n$. Write $K_n(v)$ for the connected open subgraph containing $v$, and $L_n$ for the maximum number of vertices of any open component. Diameters of connected graphs are their intrinsic shortest-path diameters; thus paths used to measure $\operatorname{diam}(K_n(v))$ must use open edges.

Suppose there is a deterministic vertex $v_n\in V(G_n)$ for every $n$ such that
\[
\mathbb P_{1/2}\!\left(\operatorname{diam}(K_n(v_n))>
\tfrac12\operatorname{diam}(G_n)\right)>\tfrac12.
\]
Must there exist $\delta>0$ such that $\mathbb P_{1/2}(L_n\ge\delta|V(G_n)|)\to1$?

This spells out the source's ``giant component with high probability'' as a positive fraction of all vertices, with one fixed fraction along the sequence. The graph family need not be vertex-transitive, and the distinguished vertex need not be typical or randomly chosen. The hypothesis gives only a probability exceeding one half of a large-diameter cluster at one vertex. The conclusion asks for a high-probability macroscopic component somewhere in the graph, at the same parameter $1/2$, without increasing the edge-retention probability.
\subsection{Short English statement}
Does a greater-than-half chance of a large-diameter cluster at one vertex of an expander force a giant cluster with high probability?
\subsection{Sources}
\begin{itemize}
\item[S1] ULAM.AI and the UnsolvedMath Contributors, supplied problems.json, record 10000044 (AMR-099-0044), AMR Open Problem Lists. Catalogue identity and source transcription; CC BY 4.0.
\url{https://huggingface.co/datasets/ulamai/UnsolvedMath}.
\item[S2] Benjamini - Coarse Geometry and Randomness (Saint-Flour notes), Open problem 11.7, PDF page 86. Original source indexed by AMR-099-0044.
\url{https://www.wisdom.weizmann.ac.il/~itai/stflouraug24.pdf#page=86}.
\item[S3] I. Benjamini, Coarse Geometry and Randomness, primary Saint-Flour notes; the numbered question and its surrounding definitions.
\url{https://arquivo.pt/noFrame/replay/20201231041548id_/http://www.wisdom.weizmann.ac.il/~itai/stflouraug24.pdf}.
\item[S4] I. Benjamini, Euclidean vs. Graph Metric, primary numbered question and conventions.
\url{https://arquivo.pt/noFrame/replay/20201231041538id_/http://www.wisdom.weizmann.ac.il/~itai/erd100.pdf}.
\end{itemize}
\end{document}
